\documentclass[11pt]{article}
\usepackage[margin=1in]{geometry}
\usepackage{amsmath,amssymb,amsthm,hyperref}
\newtheorem{theorem}{Theorem}
\newtheorem{lemma}{Lemma}
\newcommand{\res}{\operatorname{Res}}
\title{Gamma positivity of the diagonal numerators of the blade array}
\author{Principia Math In Progress}
\date{18 September 2026}
\begin{document}
\maketitle

\section*{Task 1: Complete solution of MathDB 340676}
The target is Conjecture 8.4 of Early's current preprint~\cite{Early}.
Its symmetry and coefficient-sum assertions were already proved by Donghyun
Kim in Appendix A of that paper. The remaining unimodality assertion follows
from the stronger gamma positivity proved below. We use Early's established
enumeration formula, Proposition 8.3:
\[
 T_{n,k}=\sum_{i=1}^n S(n,i)c(i-1,k-1) \qquad(1\leq k\leq n),
\]
where $S$ and $c$ are the Stirling numbers of the second and unsigned first
kinds. To fix the diagonal indexing precisely, put
\[
 H_j(x)=\sum_{r\geq0}T_{j+r+1,r+1}x^r,
 \qquad P_j(x)=(1-x)^{2j+1}H_j(x),\quad j\geq0.
\]
This gives the source's numerators $1$, $1+x$, $1+10x+x^2$, and so on.

\begin{theorem}\label{main}
For every $j\geq0$ there are nonnegative integers $q_{j,r}$ such that
\begin{equation}\label{gamma}
 P_j(x)=\sum_{r=0}^{\lfloor j/2\rfloor}
           4^r q_{j,r}x^r(1+x)^{j-2r},\qquad q_{j,0}=1.
\end{equation}
Consequently $P_j$ has symmetric, unimodal, nonnegative integer
coefficients, and
\[
 P_j(1)=\frac{(2j)!}{j!}.
\]
Its coefficients strictly increase up to the middle, apart from the
necessary pair of equal central coefficients when $j$ is odd.
\end{theorem}

All series manipulations in the proof are formal, over fields of
characteristic zero; no analytic convergence is required.

\subsection*{The row series and a formal inverse}
Set
\[
 R_m(y)=\sum_{k=0}^m T_{m+1,k+1}y^k,\qquad
 G(z)=-\log(2-e^z).
\]
The established row exponential generating function is
\begin{equation}\label{row}
 \sum_{m\geq0}R_m(y)\frac{z^m}{m!}=e^z(2-e^z)^{-y}=e^ze^{yG(z)}.
\end{equation}
This formula is also recorded by Skirrow in OEIS A320280~\cite{OEIS}.
For completeness, it follows directly from Early's formula: since
$\sum_k c(i-1,k)y^k=y^{\overline{i-1}}$ and
$\sum_{n\geq i}S(n,i)z^n/n!=(e^z-1)^i/i!$, differentiation in $z$ gives
\[
 \sum_{m\geq0}R_m(y)\frac{z^m}{m!}
 =e^z\sum_{i\geq1}y^{\overline{i-1}}
                    \frac{(e^z-1)^{i-1}}{(i-1)!}.
\]
The binomial theorem proves~\eqref{row}. In particular,
\begin{equation}\label{entry}
 T_{j+r+1,r+1}=\frac{(j+r)!}{r!}[z^{j+r}]e^zG(z)^r.
\end{equation}

Let $w=w(u,x)$ be the unique series with $w(0,x)=0$ satisfying
\begin{equation}\label{inverse}
 u=w-xG(w)=w+x\log(2-e^w).
\end{equation}
The linear coefficient $1-x$ is invertible in $\mathbb Q[[x]]$, so this
inverse exists as a formal series in $u$. We claim that
\begin{equation}\label{residue}
 H_j(x)=j![u^j]\frac{d}{du}e^{w(u,x)}.
\end{equation}
Indeed, changing variables in a formal residue and expanding in $x$ gives
\begin{align*}
 j![u^j]\frac{d}{du}e^{w(u,x)}
 &=j!\res_{z=0}\frac{e^z}{(z-xG(z))^{j+1}}\,dz\\
 &=\sum_{r\geq0}\frac{(j+r)!}{r!}x^r
             [z^{j+r}]e^zG(z)^r=H_j(x).
\end{align*}
The formal change of variables is valid for an invertible series with zero
constant term. Expansion in $x$ makes each displayed residue coefficient
well defined, and the last equality is~\eqref{entry}.

\subsection*{A differential equation with nonnegative coefficients}
Write $A=1+x$, $B=1-x$, and $V=e^{w(u,x)}$. Differentiating~\eqref{inverse},
with $V(0,x)=1$, yields
\[
 \frac{dV}{du}=\frac{V(2-V)}{2-(1+x)V}.
\]
Introduce a new variable $s$ and a parameter $h$ by
\[
 s=\frac{A}{B^2}u,\qquad
 V=1+\frac BA Y,\qquad h=\frac{4x}{A^2}.
\]
Since $A^2-B^2=4x$, the resulting initial-value problem is
\begin{equation}\label{positiveODE}
 \frac{dY}{ds}=f_h(Y),\quad Y(0,h)=0,\qquad
 f_h(y)=1+y+\frac{hy^2}{1-y}.
\end{equation}
Define, first with $h$ an independent indeterminate,
\[
 Q_j(h)=\left.\left(f_h(y)\frac{d}{dy}\right)^j f_h(y)\right|_{y=0}.
\]
Repeated differentiation of~\eqref{positiveODE} shows that
$Q_j(h)=\partial_s^{j+1}Y(0,h)$. Since
$V_u=B^{-1}Y_s$ and $s_u=A/B^2$, equation~\eqref{residue} becomes
\begin{equation}\label{Qidentity}
 H_j(x)=\frac{A^j}{B^{2j+1}}Q_j\!\left(\frac{4x}{A^2}\right).
\end{equation}

\begin{lemma}\label{positive}
$Q_j(h)\in\mathbb Z_{\geq0}[h]$, $\deg Q_j\leq\lfloor j/2\rfloor$, and
$Q_j(0)=1$.
\end{lemma}
\begin{proof}
Expand $f_h(y)=1+y+h\sum_{a\geq2}y^a$. Multiplication by $f_h$ and
differentiation in $y$ preserve nonnegative integer coefficients.
At each fixed $j$ only finitely many terms contribute at $y=0$.

For the degree bound, assign the monomial $h^r y^a$ weight $a-2r$.
Every monomial of $f_h$ has weight at least zero. Differentiation lowers
weight by one, and multiplication by $f_h$ cannot lower it. Thus every
monomial of $(f_h\,d/dy)^j f_h$ has weight at least $-j$.
After evaluation at $y=0$, a surviving $h^r$ must satisfy $-2r\geq-j$.
Finally, at $h=0$, the operator is $(1+y)d/dy$ and fixes $1+y$,
giving $Q_j(0)=1$.
\end{proof}

Writing $Q_j(h)=\sum_rq_{j,r}h^r$ in~\eqref{Qidentity} and multiplying
by $B^{2j+1}$ proves~\eqref{gamma}. Each polynomial
$x^r(1+x)^{j-2r}$ is symmetric about $j/2$ and unimodal when its
coefficient vector is padded with zeros to length $j+1$. Nonnegative sums
preserve these two properties because all summands have the same center.
The $r=0$ summand $(1+x)^j$ has coefficient one and gives strict growth
to the middle. This proves the new unimodality assertion.

For an additional derivation of the already established coefficient sum,
specialize~\eqref{positiveODE} at $h=1$. Then $Y_s=1/(1-Y)$, so
$Y=1-\sqrt{1-2s}$ and $Y_s=(1-2s)^{-1/2}$. Therefore
\[
 Q_j(1)=(2j-1)!!,\qquad
 P_j(1)=2^jQ_j(1)=\frac{(2j)!}{j!},
\]
including $j=0$ with $(-1)!!=1$. This completes the theorem.

\subsection*{Examples, verification, and scope}
The first polynomials $Q_j$ are
\[
 1,\quad1,\quad1+2h,\quad1+14h,\quad
 1+88h+16h^2,\quad1+628h+316h^2.
\]
The accompanying exact verifier compares the original Stirling sums with
the gamma formula, checks a separate row recurrence, solves the cleared
differential equation independently at six integer parameters, and checks
the source tables and coefficient properties. These finite checks support
reproducibility; the proof above treats every diagonal.

Early's enumeration formula, Kim's symmetry and factorial-sum proofs, and
Skirrow's recorded row generating function are prior inputs and are not
claimed as new. The contribution here is the positive differential-operator
formula for the diagonal gamma coefficients and the resulting all-diagonal
unimodality. Review is by the authoring agent only, without independent
peer review or formal proof-assistant certification.

\section*{Task 2: A combinatorial interpretation of the gamma coefficients}
Find a direct statistic on the canonical blade objects whose distribution
gives $q_{j,r}$, or explain the factor $4^r$ by a decoration of those objects.
Equation~\eqref{positiveODE} suggests increasing trees with $r$ vertices
having at least two children, but a correspondence with blades is still needed.

\section*{Task 3: Zeros and interlacing}
Determine whether every $Q_j(h)$ has only real negative zeros, and whether
the nonconstant successive polynomials satisfy an interlacing relation.
The proof establishes coefficient positivity, not these stronger properties.

\section*{Task 4: A limiting coefficient distribution}
For $\Pr(X_j=a)=[x^a]P_j(x)/P_j(1)$, symmetry gives
$\mathbb E X_j=j/2$. Determine the variance asymptotics and whether the
centered, normalized variables satisfy a central limit theorem.

\section*{Task 5: Refine the branching degrees}
Replace $hy^2/(1-y)$ by $\sum_{a\geq2}h_a y^a$ in~\eqref{positiveODE}.
Interpret the resulting nonnegative multivariate coefficients and identify
a refinement of the blade enumeration that specializes back to $P_j$.
These four followups are questions, not additional solved MathDB problems.

\begin{thebibliography}{9}
\bibitem{Early} N. Early, with an appendix by D. Kim,
\emph{Honeycomb Tessellations and Graded Permutohedral Blades},
\href{https://arxiv.org/abs/1810.03246v4}{arXiv:1810.03246v4},
3 February 2025, Proposition 8.3, Conjecture 8.4, and Appendix A.
\bibitem{OEIS} OEIS Foundation, \href{https://oeis.org/A320280}{A320280},
row generating functions recorded by N. L. Skirrow, 22 December 2025.
\end{thebibliography}
\end{document}
